\section{Reproducibility}
\label{app:repro}

All runs: Intel Xeon W-2295 (18C/36T), 128\,GB RAM, NVIDIA RTX A5000 (24\,GB), Windows 11,
Python 3.12, PyTorch 2.5.1+cu121, PyG 2.8.0.
% source: README.md (Hardware)
Every reported number recomputes from a run's stored metrics; an audit script re-derives 18
headline rows mechanically (18 traceable, 0 mismatched), and a frozen-numbers file pins every
headline to its source with a failing test on change.
% source: README.md; docs/REPRODUCE.md §13 (row 6); results/FROZEN.json

\paragraph{Data integrity.}
Edge identifiers are row positions after a stable sort on timestamps, so a manifest check pins
the inputs:
% source: docs/REPRODUCE.md §0

\begin{center}\small
\begin{tabular}{lll}
\toprule
File & SHA-256 (first 16) & Bytes \\
\midrule
\texttt{HI-Small\_Trans.csv} & \texttt{b19d39f515523373} & 475{,}664{,}283 \\ % source: docs/REPRODUCE.md §0
\texttt{HI-Small\_accounts.csv} & \texttt{786808526e33cfc4} & 34{,}053{,}187 \\ % source: docs/REPRODUCE.md §0
\texttt{HI-Small\_Patterns.txt} & \texttt{2c546b5ce6009e73} & 323{,}844 \\ % source: docs/REPRODUCE.md §0
\texttt{formatted\_transactions.csv} & \texttt{3b9678458eeece7e} & 266{,}605{,}148 \\ % source: docs/REPRODUCE.md §0
\bottomrule
\end{tabular}
\end{center}

\paragraph{End-to-end check.}
With nothing retrained, the data manifest passes, the research test suite passes 258 tests,
and the ring-building command regenerates the released rings byte-identically (108 rings in
62\,s).
% source: docs/REPRODUCE.md §13 (rows 1, 2a, 7)
Measured training times: detector about 15\,h for 5 seeds; cross-scoring about 8\,h; MEGA-GNN
about 18\,h; statistics about 45\,min.
% source: docs/REPRODUCE.md §2, §4, §7, §8

\section{Additional results}
\label{app:tables}

\paragraph{Scope conventions.}
Tables~\ref{tab:scope} and~\ref{tab:component-scope} give the numbers behind
Figure~\ref{fig:scope}.

% Source: results/tables/r14_scope_sensitivity.csv, as summarised in
%         docs/paper/evaluation_protocol.md §9 (same shipped runs re-scored; no model re-run)
\begin{table}[t]
\centering
\caption{The scope convention moves ring-F1 about five times more than the entire method does.
The same shipped runs re-scored under \texttt{full\_membership} (a ring's member set is all its
accounts) against the reported \texttt{test\_visible}; the method ordering survives under both.}
\label{tab:scope}
\small
\begin{tabular}{lccc}
\toprule
Detector / method & \texttt{test\_visible} (reported) & \texttt{full\_membership} & Gap \\
\midrule
GINe B0 & 0.5213 & 0.1437 & $-0.378$ \\ % source: r14_scope_sensitivity.csv (gine_causal b0 rows, ring_f1)
GINe B1b & 0.5553 & 0.2073 & $-0.348$ \\ % source: r14_scope_sensitivity.csv (gine_causal b1b rows)
GINe proposed & 0.5819 & 0.2297 & $-0.352$ \\ % source: r14_scope_sensitivity.csv (gine_causal proposed rows)
GINe full & 0.5894 & 0.2428 & $-0.347$ \\ % source: r14_scope_sensitivity.csv (gine_causal full rows)
MEGA B0 & 0.5208 & 0.1595 & $-0.361$ \\ % source: r14_scope_sensitivity.csv (mega b0 rows)
MEGA full & 0.6074 & 0.2252 & $-0.382$ \\ % source: r14_scope_sensitivity.csv (mega full rows)
\bottomrule
\end{tabular}
\end{table}

% Source: results/tables/v36_component_scope_wci.csv (v36_component_scope.csv re-scored with the
%         weighted ring bootstrap by evaluation/rescore_weighted_intervals.py; point estimates identical)
\begin{table}[t]
\centering
\caption{A third scope convention: each in-window connected piece of a truth ring scored as
its own target (\texttt{component}: 254 targets), or only pieces of $\ge 3$ accounts
(\texttt{component\_ge3}: 108 targets in 97 rings). Ring-F1 with 95\,\% bootstrap CIs over
targets; same shipped runs, nothing re-selected. Denominators differ between conventions, so
the right-hand columns are never headline numbers.}
\label{tab:component-scope}
\small
\begin{tabular}{llccc}
\toprule
Detector & Method & \texttt{test\_visible} & \texttt{component} & \texttt{component\_ge3} \\
\midrule
GINe & B0 & \ci{0.521}{0.429}{0.607} & \ci{0.444}{0.371}{0.514} & \ci{0.673}{0.576}{0.761} \\ % source: v36_component_scope_wci.csv (gine_causal, b0)
GINe & B1b & \ci{0.555}{0.470}{0.632} & \ci{0.488}{0.412}{0.562} & \ci{0.700}{0.628}{0.767} \\ % source: v36_component_scope_wci.csv (gine_causal, b1b)
GINe & proposed & \ci{0.582}{0.490}{0.668} & \ci{0.504}{0.425}{0.572} & \ci{0.766}{0.678}{0.840} \\ % source: v36_component_scope_wci.csv (gine_causal, proposed)
GINe & full & \ci{0.589}{0.502}{0.675} & \ci{0.515}{0.438}{0.583} & \ci{0.740}{0.656}{0.812} \\ % source: v36_component_scope_wci.csv (gine_causal, full)
MEGA & B0 & \ci{0.521}{0.432}{0.609} & \ci{0.441}{0.366}{0.510} & \ci{0.674}{0.588}{0.762} \\ % source: v36_component_scope_wci.csv (mega, b0)
MEGA & proposed & \ci{0.598}{0.506}{0.683} & \ci{0.517}{0.437}{0.583} & \ci{0.777}{0.707}{0.846} \\ % source: v36_component_scope_wci.csv (mega, proposed)
MEGA & full & \ci{0.607}{0.521}{0.688} & \ci{0.530}{0.452}{0.595} & \ci{0.774}{0.709}{0.837} \\ % source: v36_component_scope_wci.csv (mega, full)
\bottomrule
\end{tabular}
\end{table}


\paragraph{Validation-to-test gap.}
Table~\ref{tab:val-test-gap} gives the per-seed edge F1 behind the scoring-batch footnote in
Section~\ref{sec:results-swap}.

% Source: results/notes/p7_val_test_gap.csv (Multi-GINe), results/notes/p7_val_test_gap_sb2048.csv (MEGA-GNN, scoring batch 2048),
% shipped batch-8192 values from results/tables/r6_detector_swap.csv (edge_f1_per_seed)
\begin{table}[t]
\centering
\caption{Validation-to-test edge-F1 gap per seed. The threshold is chosen on validation
(maximum F1 over 400 validation-score quantiles) and applied to test. MEGA-GNN rows use
scoring batch 2048, as everywhere in the paper; the last two columns give the test F1 at the
original scoring batch of 8192 and the difference. Every MEGA-GNN seed scores higher on test
than on validation; Multi-GINe shows no consistent gap.}
\label{tab:val-test-gap}
\small
\begin{tabular}{lcccccc}
\toprule
Detector & Seed & Val F1 & Test F1 & Test $-$ val & Test F1 (batch 8192) & 2048 $-$ 8192 \\
\midrule
Multi-GINe & 1 & 0.6140 & 0.6472 & $+0.0332$ & -- & -- \\ % source: results/notes/p7_val_test_gap.csv
Multi-GINe & 2 & 0.5609 & 0.6065 & $+0.0456$ & -- & -- \\ % source: results/notes/p7_val_test_gap.csv
Multi-GINe & 3 & 0.6275 & 0.6575 & $+0.0300$ & -- & -- \\ % source: results/notes/p7_val_test_gap.csv
Multi-GINe & 4 & 0.6140 & 0.5870 & $-0.0270$ & -- & -- \\ % source: results/notes/p7_val_test_gap.csv
Multi-GINe & 5 & 0.5508 & 0.4806 & $-0.0702$ & -- & -- \\ % source: results/notes/p7_val_test_gap.csv
\midrule
MEGA-GNN & 1 & 0.5564 & 0.6309 & $+0.0745$ & 0.6435 & $-0.0126$ \\ % source: results/notes/p7_val_test_gap_sb2048.csv; results/tables/r6_detector_swap.csv
MEGA-GNN & 2 & 0.5892 & 0.6784 & $+0.0893$ & 0.6813 & $-0.0029$ \\ % source: results/notes/p7_val_test_gap_sb2048.csv; results/tables/r6_detector_swap.csv
MEGA-GNN & 3 & 0.6050 & 0.6658 & $+0.0608$ & 0.6671 & $-0.0013$ \\ % source: results/notes/p7_val_test_gap_sb2048.csv; results/tables/r6_detector_swap.csv
MEGA-GNN & 4 & 0.6558 & 0.7512 & $+0.0955$ & 0.7463 & $+0.0049$ \\ % source: results/notes/p7_val_test_gap_sb2048.csv; results/tables/r6_detector_swap.csv
MEGA-GNN & 5 & 0.6095 & 0.7014 & $+0.0919$ & 0.7083 & $-0.0069$ \\ % source: results/notes/p7_val_test_gap_sb2048.csv; results/tables/r6_detector_swap.csv
\midrule
Multi-GINe & mean & 0.5935 & 0.5957 & $+0.0023$ & -- & -- \\
MEGA-GNN & mean & 0.6032 & 0.6856 & $+0.0824$ & 0.6893 & $-0.0037$ \\
\bottomrule
\end{tabular}
\end{table}


\paragraph{Seed budget and windows.}
The validation-chosen budget is not test-optimal: the full pipeline at 0.15\,\% scores
$0.606 \pm 0.033$ against \ci{0.589}{0.502}{0.675} at the validation-chosen 0.1\,\%, a gap
inside seed noise, and we do not move the operating point (Table~\ref{tab:robustness}).
Ring-level numbers are stable between the primary and whole windows
(Table~\ref{tab:windows}).
% source: docs/FINDINGS.md (section E; section D)

% Source: docs/FINDINGS.md §E (aggregated from results/tables/r9_robustness.csv, per-seed rows)
\begin{table}[t]
\centering
\caption{Robustness of the quiet-member result to the seed budget $k$ (primary window, 5
seeds). $k^{*}=0.1\,\%$ was chosen on validation and is not changed; non-$k^{*}$ rows are
mean $\pm$ std over seeds (a robustness check, not a headline). Added-member precision here was measured with day-folded expansion-policy inputs; with leak-free (ring-folded) inputs the reference full pipeline's value is 0.817 (Table~\ref{tab:precision-paired}).}
\label{tab:robustness}
\small
\begin{tabular}{llccc}
\toprule
$k$ (seed edges) & Method & Ring F1 & Quiet-member recall & Added-member precision \\
\midrule
0.05\,\% (432) & proposed & $0.433 \pm 0.069$ & $0.140 \pm 0.041$ & $0.970 \pm 0.016$ \\ % source: docs/FINDINGS.md (E1)
0.05\,\% & full & $0.450 \pm 0.066$ & $0.140 \pm 0.041$ & $0.970 \pm 0.016$ \\ % source: docs/FINDINGS.md (E2)
\textbf{0.10\,\% (864), shipped $k^{*}$} & proposed & \ci{0.582}{0.490}{0.668} & \ci{0.209}{0.142}{0.290} & \ci{0.926}{0.899}{0.951} \\ % source: docs/FINDINGS.md (E3; quiet recall as the 0.2092 point estimate per the 2026-11-12 note in §A)
0.10\,\%, shipped $k^{*}$ & full & \ci{0.589}{0.502}{0.675} & \ci{0.203}{0.136}{0.284} & \ci{0.926}{0.899}{0.948} \\ % source: docs/FINDINGS.md (E4)
0.15\,\% (1{,}296) & proposed & $0.593 \pm 0.044$ & $0.235 \pm 0.020$ & $0.783 \pm 0.069$ \\ % source: docs/FINDINGS.md (E5)
0.15\,\% & full & $0.606 \pm 0.033$ & $0.232 \pm 0.015$ & $0.782 \pm 0.069$ \\ % source: docs/FINDINGS.md (E6)
\bottomrule
\end{tabular}
\end{table}

% Source: docs/FINDINGS.md §D (D1–D4), from results/tables/r9_ring_ci.csv
\begin{table}[t]
\centering
\caption{The headline rows on both windows. Ring-F1 is stable across the two windows
($\pm 0.01$ for every method except B0); the tail's 59\,\% base rate inflates edge-level
numbers, not these ring-level ones.}
\label{tab:windows}
\small
\begin{tabular}{lcccc}
\toprule
 & \multicolumn{2}{c}{Ring F1} & \multicolumn{2}{c}{Quiet-member recall} \\
\cmidrule(lr){2-3}\cmidrule(lr){4-5}
Method & primary (d.\,8--9) & whole (d.\,8--17) & primary & whole \\
\midrule
B0 & \ci{0.521}{0.429}{0.607} & \ci{0.471}{0.395}{0.555} & \ci{0.0052}{0.000}{0.015} & \ci{0.0000}{0.000}{0.000} \\ % source: docs/FINDINGS.md (D1, point estimate corrected 2026-11-12); r9_quiet_members.csv (b0,primary,ALL)
B1b (rule) & \ci{0.555}{0.470}{0.632} & \ci{0.552}{0.475}{0.634} & \ci{0.223}{0.163}{0.299} & \ci{0.252}{0.192}{0.319} \\ % source: docs/FINDINGS.md (D2)
Proposed (policy) & \ci{0.582}{0.490}{0.668} & \ci{0.577}{0.490}{0.667} & \ci{0.209}{0.142}{0.290} & \ci{0.237}{0.171}{0.306} \\ % source: docs/FINDINGS.md (D3: 0.2092, corrected 2026-11-12)
Full & \ci{0.589}{0.502}{0.675} & \ci{0.587}{0.499}{0.674} & \ci{0.203}{0.136}{0.284} & \ci{0.227}{0.165}{0.291} \\ % source: docs/FINDINGS.md (D4)
\bottomrule
\end{tabular}
\end{table}


\paragraph{Pattern typing.}
The typing classifier is correct on \ci{0.599}{0.549}{0.648} of 401 matched rings, against
\ci{0.551}{0.499}{0.601} for a rule, with overlapping intervals; accuracy does not change with
the window (Table~\ref{tab:typing}). RANDOM (0.11) and BIPARTITE (0 of 5) remain weak.
% source: results/tables/r5_typing_ci.csv; docs/FINDINGS.md (B10)

% Source: results/tables/r5_typing_ci.csv
\begin{table}[t]
\centering
\caption{Typing accuracy over matched rings (a per-ring classification rate, not a ring-F1;
bootstrap CIs over matched rings). Typing does not change with the window.}
\label{tab:typing}
\small
\begin{tabular}{lccc}
\toprule
Window & Matched rings & Head accuracy [CI] & Rule accuracy [CI] \\
\midrule
primary (days 8--9) & 401 & \ci{0.599}{0.549}{0.648} & \ci{0.551}{0.499}{0.601} \\ % source: r5_typing_ci.csv (primary rows)
whole (days 8--17) & 433 & \ci{0.594}{0.547}{0.640} & \ci{0.547}{0.499}{0.591} \\ % source: r5_typing_ci.csv (whole rows)
\bottomrule
\end{tabular}
\end{table}


\paragraph{Held-out pattern types.}
We removed one pattern type from every label set the ring layer trains on, in train and
validation together. Holding out BIPARTITE removes none of the policy's 244 positive training
rows, because BIPARTITE quiet edges never share an endpoint with a seed of their own ring; the
pattern is unreachable by endpoint-adjacent expansion whether or not it is held out, and its
unchanged number is not evidence of transfer. Holding out CYCLE removes 14.8\,\% of the
positives --- every positive for that topology --- and the policy still adds CYCLE members at
undiminished precision (Table~\ref{tab:heldout-cycle}); the recall intervals are wide and
cannot exclude a loss of up to 6.5 points.
% source: docs/paper/generalisation.md §2; docs/FINDINGS.md (B6)

% Source: docs/paper/generalisation.md §2 (CYCLE arm), from results/tables/r7_heldout.csv
\begin{table}[t]
\centering
\caption{Held-out CYCLE: the expansion policy trained with every
CYCLE positive removed, evaluated on the 22 CYCLE truth rings (paired bootstrap).}
\label{tab:heldout-cycle}
\small
\begin{tabular}{lcccc}
\toprule
CYCLE rings & All patterns & CYCLE held out & $\Delta$ [95\,\% CI] & $p$ \\
\midrule
Quiet-member recall & 0.1614 & 0.1344 & $+0.0270$ [$-0.0062$, $+0.0647$] & 0.130 \\ % source: docs/paper/generalisation.md §2 (CYCLE table)
Ring recall & 0.3818 & 0.3455 & $+0.0364$ [$-0.0125$, $+0.0728$] & 0.286 \\ % source: docs/paper/generalisation.md §2
Added-member precision & 0.952 (40/42) & 0.970 (32/33) & --- & --- \\ % source: docs/paper/generalisation.md §2; r7_heldout.csv (l3_cycle,full,CYCLE)
\bottomrule
\end{tabular}
\end{table}


\paragraph{Per-pattern detection and time-to-detection.}
Table~\ref{tab:per-pattern} and Figure~\ref{fig:appendix-figs} break results down by pattern.
Time-to-detection is uninformative on this dataset --- median 0\,h, because a ring's first
in-window transaction is usually already a seed --- and is not reported further.
% source: docs/report.md §7.2; results/tables/r9_ttd.csv

% Source: docs/report.md §7.2 "Detection by ring type" (full pipeline, 5 seeds summed;
%         computed by evaluation/unmatched_rings.py)
\begin{table}[t]
\centering
\caption{Detection by pattern type (full pipeline, 5 seeds summed, primary window). ``Found
\%'' excludes the below-\texttt{min\_size} column from its denominator.}
\label{tab:per-pattern}
\small
\begin{tabular}{lrrrrr}
\toprule
Pattern & Found & $<3$ visible, unmatched & Never flagged & Merged/partial & Found \% \\
\midrule
SCATTER-GATHER & 78 & 25 & 5 & 2 & 92\,\% \\ % source: docs/report.md §7.2 (per-type table)
FAN-OUT & 64 & 24 & 6 & 6 & 84\,\% \\ % source: docs/report.md §7.2
GATHER-SCATTER & 117 & 43 & 17 & 8 & 82\,\% \\ % source: docs/report.md §7.2
FAN-IN & 56 & 21 & 12 & 1 & 81\,\% \\ % source: docs/report.md §7.2
CYCLE & 42 & 40 & 18 & 10 & 60\,\% \\ % source: docs/report.md §7.2
RANDOM & 35 & 25 & 21 & 4 & 58\,\% \\ % source: docs/report.md §7.2
BIPARTITE & 5 & 25 & 41 & 9 & 9\,\% \\ % source: docs/report.md §7.2
STACK & 4 & 40 & 46 & 20 & 6\,\% \\ % source: docs/report.md §7.2
\bottomrule
\end{tabular}
\end{table}


\begin{figure}[t]
\centering
\begin{subfigure}[t]{0.48\textwidth}
  \includegraphics[width=\linewidth]{f2_per_pattern_recall.png}
  \caption{Per-pattern recall.}
  % source: results/figures/f2_per_pattern_recall.png
\end{subfigure}\hfill
\begin{subfigure}[t]{0.48\textwidth}
  \includegraphics[width=\linewidth]{r4_val_ringf1_vs_seedfrac.png}
  \caption{Validation ring-F1 against seed fraction.}
  % source: results/figures/r4_val_ringf1_vs_seedfrac.png
\end{subfigure}
\caption{Supplementary figures from the repository.}
\label{fig:appendix-figs}
\end{figure}

\section{Corrections made during the project}
\label{app:corrections}

Four claims were corrected before this draft and appear in the paper only in corrected form.
(i) An earlier bound on ring recall derived from \texttt{min\_size} was circular and is
withdrawn (Section~\ref{sec:results-ceiling}). (ii) An earlier, much lower typing accuracy
came from an evaluation lookup that typed a different ring than the matched one; it was found
by an independent check from the explanation layer and fixed with a regression test. (iii) The
partially recovered rings were first read as chains the policy declined to follow; the
connectivity measurement shows most missing members were never reachable. (iv) An earlier
comparison suggested a per-account seed budget lowers baseline ring-F1; at an equal budget it
does not.
% source: docs/FINDINGS.md (B7/B9, B10, B12/B13, and the dated correction to B14)

\section{System context: explanations and platform}
\label{app:platform}

The rings feed two components that are delivered capabilities, not validated research claims.
Their numbers come from their own reports and are outside the findings file this paper is
otherwise restricted to.
% source: docs/FINDINGS.md (scope note of 2026-11-12); README.md

\paragraph{Explanations.}
Each ring receives transaction weights (calibrated score and flow share), account roles,
TreeSHAP feature attributions and a fixed-template narrative; there is no language model.
Median latency is 393\,ms per ring. On the standard fidelity metric the ranking is
indistinguishable from score-only because most matched rings contain only true transactions;
on the 52 instances that mix true and false transactions, score-only is better than the mixed
ranking ($+0.058$ [$+0.010$, $+0.133$]). An internal three-person review of 30 explanations
found the named accounts right at a 0.91 yes-rate and the stated shape plausible at 0.53; it
was not an analyst study, and an external one is designed but not run.
% source: explainability/report/xai_chapter.md §5.2, §7; docs/deep-analysis.md §11, §14

\paragraph{Platform.}
A multi-bank analyst dashboard serves the rings with per-bank isolation enforced by database
row-level security, records verdicts, and gates retraining on a frozen holdout
(Figure~\ref{fig:platform}).
% source: docs/report.md §3.1–3.2
A pooled model applied to one bank's view recovers \ci{0.088}{0.031}{0.161} of tenant-ring
pairs against \ci{0.359}{0.239}{0.500} with the pooled view (Table~\ref{tab:r10}); 172 of 174
rings span more than one bank. Pooling the five studied banks closes 11\,\% of that gap
(0.121 to 0.147).
% source: results/tables/r10_registered.json; results/tables/r10_consortium.json; docs/report.md §7.2
A drift-monitoring experiment could not be made testable on the two-day primary window: the
detector's scores have a tie plateau across the 99th percentile, and day 9 already differs
from day 8 by more than the synthetic shift.
% source: docs/R12.md (second registration); results/platform/r12b/summary.json

% Source: results/tables/r10_registered.md / r10_registered.json, as summarised in docs/report.md §7.2
% R10 is a platform-side experiment registered in docs/TENANT_SIZES.md; it is NOT in docs/FINDINGS.md.
\begin{table}[t]
\centering
\caption{Ring recall of a pooled model applied to one
bank's view vs.\ the pooled view. Banks 1, 20, 11, 118, 12; 68 tenant-ring pairs; 10{,}000
bootstrap resamples. The gap includes distribution shift, so it is not only the cost of seeing
fewer edges.}
\label{tab:r10}
\small
\begin{tabular}{lccc}
\toprule
Row & Bank alone & Pooled & Pooled $-$ bank \\
\midrule
Primary: bank's own scores, features, embeddings & \ci{0.088}{0.031}{0.161} & \ci{0.359}{0.239}{0.500} & \ci{0.271}{0.154}{0.407} \\ % source: r10_registered.json (rows.primary); docs/report.md §7.2
Pooled scores on the bank's edges & \ci{0.121}{0.054}{0.205} & 0.359 & \ci{0.238}{0.135}{0.363} \\ % source: docs/report.md §7.2 (R10 table)
Matched seed budget & \ci{0.076}{0.020}{0.151} & 0.359 & \ci{0.282}{0.167}{0.417} \\ % source: docs/report.md §7.2 (R10 table)
\bottomrule
\end{tabular}
\end{table}


\begin{figure}[t]
\centering
\begin{subfigure}[t]{0.48\textwidth}
  \includegraphics[width=\linewidth]{02-ring-queue-light.png}
  \caption{The analyst's ring queue.}
  % source: docs/screenshots/02-ring-queue-light.png
\end{subfigure}\hfill
\begin{subfigure}[t]{0.48\textwidth}
  \includegraphics[width=\linewidth]{05-ring-detail-narrative-light.png}
  \caption{A ring's narrative.}
  % source: docs/screenshots/05-ring-detail-narrative-light.png
\end{subfigure}
\caption{The multi-bank analyst platform.}
\label{fig:platform}
\end{figure}
